\chapter{Who holds the keys}

Custody is who holds the private keys. That person, or
that company, can spend. Everything else is branding.

\section{Custodial}

A custodial wallet is an account. An exchange or a
wallet company holds the keys. You log in. They can
reset the password. They can also freeze you, get
hacked, go insolvent, or hand the keys to someone with a
warrant.

Use this when you are trading on that venue, or when you
accept counterparty risk for the recovery convenience.
Do not keep long-term savings there because the app is
pretty.

\section{Non-custodial}

A non-custodial wallet is a key. You hold it. Nobody can
reset it. Nobody can help you if you lose it. dApps
talk to it directly.

This book is for that case. If someone else holds the
keys, their security program is the book, not this one.

\section{The trade}

\begin{itemize}
  \item Custodial: recovery exists. So does a third
  party who can lose or seize the funds.
  \item Non-custodial: no counterparty. You are the
  backup.
\end{itemize}

Account abstraction and social recovery try to sit
between those poles. They do not delete the question.
Someone still has to be able to spend, and you still
have to know who.

\section{What this book assumes}

You will self-custody something. Maybe a small hot
wallet. Maybe the savings. Maybe a signer seat. The
habit that scales is the same: unique secrets, careful
signing, offline backups. Start there even if the
balance is not ``serious'' yet. Serious arrives faster
than the backup does.
